\section{Pinned frontier and purpose}
The integration base for this continuation is
\texttt{KokunoYumeto/erdos-straus-foundation} commit
\texttt{6c064a07298cdba38352998364ce6b7cb3efa973}, which incorporates the
square-zero/Pell tranche on top of the earlier mixed-fibre and middle-spectrum
revision. Its complete source and handoff were read, including the obstruction
at $p=67369$ to repairing every trace without increasing the distinguished
denominator. The retained ten-stage plan still requires an original occupied
coefficient at every prescribed prime.

We prove a stronger constructive coefficient criterion that uses directions of
\emph{different orders}, including nonunits, in the original fine-defect group.
It is a cyclic subgroup-transversal construction with exact original multiplicity,
carry, and integer-section data. The cyclic calculation is proved from first principles below. No historical-priority
claim is made for either the abstract normal form or its arithmetic application. A second construction changes the shell and the word,
and gives a one-step upward return at the predecessor's decisive $67369$ example.
Neither construction assumes that its source exists for every prime.

\section{The original bounded source and its fine defect}
Fix a prime $p\equiv1\pmod4$, $p/4<a<p/2$, and $R=4a-p$. For every
\emph{original} divisor $u\mid a^2$, put
\[
 a=\prod_iq_i^{e_i},\quad u=\prod_iq_i^{f_i},\quad
 0\le f_i\le2e_i,\quad\beta_i=f_i-e_i,
\]
\[
 g=\gcd(a,u),\quad h=g^2/u,\quad r=u/g,\quad s=a/g.
\]
Then $a=hrs$, $u=hr^2$, and $\gcd(r,s)=1$. The ordered rational returns are
\begin{align}
 E:&\quad(a,hs\kappa,phr\kappa),&\kappa&=(pr+s)/R,\label{eq:E}\\
 M:&\quad(a,phs\lambda,phr\lambda),&\lambda&=(r+s)/R.\label{eq:M}
\end{align}
Direct common-denominator multiplication gives $4/p$ in each case. Their
integer gates are $R\mid4u+1$ and $R\mid p+4u$, respectively. The ordered
inverses are $u=a^2/(Ry-pa)$ in E and $u=pa^2/(Ry-pa)$ in M.
These are the inherited Yamamoto--Rosati/Type-I--II coordinates
\citep[p.~38, Lemma~2, equation~(4)]{Yamamoto}\citep[\S2, Propositions~2.2
and~2.6]{ET}, not a new parametrization.
Lopez's structural classification supplies additional literature context for
the difficult class $p\equiv1\pmod{24}$ and for Type~II solution forms
\citep[Introduction and the Type~A/B definitions]{Lopez}; no theorem of that
paper is used to prove the new subgroup-chain or factor-exchange results below.

The rational tail sums are $(a+pu)^2/(Ru)$ and $p(a+u)^2/(Ru)$.
Consequently a trace, meaning an integral tail sum, is exactly the square gate
$R\mid G^2$, where $G=4u+1$ or $G=p+4u$. Its common tail denominator is
$R/\gcd(R,G)$. Define
\[
 K=\prod_{\ell^v\parallel R}\ell^{\lceil v/2\rceil},\quad
 D=R/K,\quad\delta=R/D^2.
\]
Then $R=KD=\delta D^2$, $D\mid K$, and the ideal $K\mathbb Z/R\mathbb Z$
is square-zero. Here $\delta$ is the positive squarefree part, not the radical.
For $\eta=1$ in E and $\eta=0$ in M, retain the original residue ratio
\[
 z=p^\eta u/a=p^\eta\prod_iq_i^{\beta_i}\pmod R.
\]
Trace is exactly $z\equiv-1\pmod K$. Its unique fine coordinate is
\begin{equation}\label{eq:defect}
 z=-1+Kc,\qquad c\in\mathbb Z/D.
\end{equation}
The full gate is $c=0$; the original denominator is $D/\gcd(D,c)$.
No integer coefficient is reduced modulo $D$.

For completeness, partition each original interval $[-e_i,e_i]$ by residue
modulo $o_i=\operatorname{ord}_K(q_i)$. Each nonempty part has the literal form
\[
 \beta_i=b_i+o_it_i,\qquad0\le t_i<N_i.
\]
A product part is a trace part when its base is $-1\pmod K$. Put
\[
 \lambda_i=(q_i^{o_i}-1)/K\pmod D,
\]
where the division is evaluated modulo $R$ before passing to $D$. Since
$(1+K\lambda)^t=1+tK\lambda\pmod R$, its fine coordinate is
$c_0-\sum_i\lambda_it_i$. Thus the complete original count in that part is
\begin{equation}\label{eq:C}
 C_{c_0},\qquad
 \mathcal C=\prod_i\left([0]+[\lambda_i]+\cdots+[(N_i-1)\lambda_i]\right)
 \in\mathbb Z[\mathbb Z/D].
\end{equation}
The inverse reads $(b_i,t_i)$ from $\beta_i$ and reconstructs
$u=\prod_iq_i^{e_i+\beta_i}$. Every vector in the full exponent box belongs to
one product part. An explicitly specified subrectangle of such a part also
has formula~\eqref{eq:C}; its omitted words remain in the complete source.
These source identities are inherited from the supplied square-zero tranche
\citep[\S\S2--4]{Previous}; the stronger coefficient criterion follows next.

\section{The exact affine relation between the two channels}
For the same divisor $u$, put
\[
 G_E(u)=4u+1,\qquad G_M(u)=p+4u.
\]
The preceding trace calculation gives
\[
 E\text{ trace}\Longleftrightarrow K\mid G_E(u),\qquad
 M\text{ trace}\Longleftrightarrow K\mid G_M(u).
\]

\begin{theorem}[Common traces and the two-channel lattice]
\label{thm:affine-channels}
The same original divisor $u\mid a^2$ is a trace in both channels if and only if
\begin{equation}\label{eq:common-trace}
 K\mid(4u+1)\quad\text{and}\quad K\mid(p-1).
\end{equation}
Equivalently, one channel is a trace and $p\equiv1\pmod K$. More strongly, if
\[
 \mathcal T_E=\{u\mid a^2:K\mid4u+1\},\qquad
 \mathcal T_M=\{u\mid a^2:K\mid p+4u\},
\]
then these two trace sets are disjoint when $K\nmid p-1$ and equal when
$K\mid p-1$.

On the common-trace locus the exact fine coordinates are
\begin{equation}\label{eq:channel-coordinates}
 c_E\equiv\frac{4u+1}{K}\pmod D,\qquad
 c_M\equiv p^{-1}\frac{p+4u}{K}\pmod D.
\end{equation}
Writing
\[
 \omega=(p-1)/K,\qquad \bar\omega=\omega\pmod D,
\]
one has the oriented affine relation
\begin{equation}\label{eq:affine-shift}
 c_E=c_M-\bar\omega\quad\text{in }\mathbb Z/D.
\end{equation}
Thus $M$ is full exactly at $c_M=0$, whereas $E$ is full exactly at
$c_M=\bar\omega$. The same divisor is full in both channels exactly when
$R\mid4u+1$ and $R\mid p-1$.

Let a common trace part have $M$ base coordinate $c_0$ and coefficient vector
$C$ from~\eqref{eq:C}. Its labelled full-state counts are
\begin{equation}\label{eq:paired-count}
 N_M=C_{c_0},\qquad N_E=C_{c_0-\bar\omega},\qquad
 N_{E\sqcup M}=C_{c_0}+C_{c_0-\bar\omega}.
\end{equation}
On the full coefficient lattice define
\[
 (T_{\bar\omega}f)_x=f_{x-\bar\omega},\qquad
 J=(I+T_{\bar\omega})C.
\]
Put $d=\gcd(D,\omega)$ and $L=D/d$. Since $D$ is odd, every
$\bar\omega$-orbit has odd length $L$, and
\begin{equation}\label{eq:alternating-inverse}
 2C_x=\sum_{j=0}^{L-1}(-1)^jJ_{x-j\bar\omega}.
\end{equation}
If $\mathcal O$ is the set of these $d$ orbits and
$\pi(f)_O=\sum_{x\in O}f_x\pmod2$, then
\begin{equation}\label{eq:channel-exact-sequence}
0\longrightarrow\mathbb Z^{\mathbb Z/D}
 \xrightarrow{I+T_{\bar\omega}}\mathbb Z^{\mathbb Z/D}
 \xrightarrow{\pi}(\mathbb Z/2)^d\longrightarrow0
\end{equation}
is exact. In particular the two-channel operator is injective, and its only
integral cokernel is the displayed orbit-parity data.

For $C_x\ge0$, uniform labelled two-channel positivity is equivalent to
\begin{equation}\label{eq:support-cover}
 \operatorname{supp}(C)\cup
 (\operatorname{supp}(C)+\bar\omega)=\mathbb Z/D.
\end{equation}
It forces the sharp necessary bound
\begin{equation}\label{eq:support-bound}
 |\operatorname{supp}(C)|\ge \frac{D+d}{2}.
\end{equation}
The adjacency condition~\eqref{eq:support-cover}, not cardinality alone, is
the exact criterion.
\end{theorem}

\begin{proof}
The difference $G_M(u)-G_E(u)=p-1$ proves both the common-trace criterion and
the trace-set dichotomy. Since $4a\equiv p\pmod R$ and $a,p$ are units modulo
$R$,
\[
 pu\,a^{-1}+1\equiv4u+1\pmod R,
 \qquad
 u\,a^{-1}+1\equiv p^{-1}(p+4u)\pmod R.
\]
Both right sides lie in $K\mathbb Z/R$ on their respective trace loci.
The map $Kc\mapsto c$ identifies $K\mathbb Z/R$ with $\mathbb Z/D$, which
proves~\eqref{eq:channel-coordinates} without discarding a factor. On the
common locus $p=1+K\omega$ and $p\equiv1\pmod D$, while
\[
 \frac{p+4u}{K}=\frac{4u+1}{K}+\omega.
\]
This proves~\eqref{eq:affine-shift}. Replacing divisibility by $K$ with
divisibility by $KD=R$ proves the full-gate statements.

For every bounded exponent vector, multiplication of
$u/a=-1+Kc_M$ by $p=1+K\omega$ gives
\[
 p(-1+Kc_M)=-1+K(c_M-\omega)+K^2\omega c_M.
\]
Because $R\mid K^2$, the last summand vanishes modulo $R$. The shift is
therefore constant over the whole common trace part; its step coefficients
and literal exponent intervals are unchanged. Formula~\eqref{eq:paired-count}
follows from the original coefficient definition~\eqref{eq:C}.

On an orbit of $T=T_{\bar\omega}$, $T^L=I$ and $L$ is odd. Hence
\[
 (I+T)(I-T+T^2-\cdots+T^{L-1})=I+T^L=2I,
\]
which proves~\eqref{eq:alternating-inverse} and injectivity. Every vector in
the image has even coordinate sum on each orbit. Conversely, if those sums
are even, every alternating numerator in~\eqref{eq:alternating-inverse} is
even, so the reconstructed $C$ is integral. The parity map is onto by taking
one standard basis vector on each orbit. This proves the exact sequence.

Finally, $J_x>0$ exactly when at least one of $C_x,C_{x-\bar\omega}$ is
positive, proving~\eqref{eq:support-cover}. Its complement has no two
consecutive vertices on any $\bar\omega$-cycle. An independent set in an odd
cycle of length $L$ has at most $(L-1)/2$ vertices, so the support has at least
$(L+1)/2$ vertices on each of $d$ cycles. This gives~\eqref{eq:support-bound},
and alternating occupancy with one forced adjacent occupied pair on each odd
cycle attains the bound.
\end{proof}

\begin{figure}[ht]
\centering
\begin{tikzpicture}[
  font=\small,
  box/.style={draw,rounded corners,minimum height=10mm,minimum width=39mm,align=center},
  map/.style={-{Latex[length=2mm]},thick}
]
\node[box,fill=blue!7] (m) {$M$ fine coordinate\\$c_M=c$};
\node[box,fill=green!9,right=25mm of m] (e) {$E$ fine coordinate\\$c_E=c-\bar\omega$};
\draw[map] (m) -- node[above] {$p=1+K\omega$} (e);
\node[below=5mm of m,align=center] {$M$ full: $c=0$};
\node[below=5mm of e,align=center] {$E$ full: $c=\bar\omega$};

\node[box,below=20mm of m,fill=black!4] (c)
  {$C\in\mathbb Z^{\mathbb Z/D}$\\original coefficient vector};
\node[box,right=25mm of c,fill=black!4] (j)
  {$J=(I+T_{\bar\omega})C$\\$J_x=C_x+C_{x-\bar\omega}$};
\draw[map] (c) -- node[above] {$I+T_{\bar\omega}$} (j);

\node[draw,rounded corners,below=10mm of $(c)!0.5!(j)$,align=center,
      text width=116mm,inner sep=5pt] (inverse)
  {$d=\gcd(D,\omega)$, $L=D/d$ is odd, and on every translation orbit\\[2pt]
   $\displaystyle 2C_x=\sum_{r=0}^{L-1}(-1)^rJ_{x-r\bar\omega}$,\qquad
   $\displaystyle J_x>0\ (\forall x)\Longleftrightarrow
   \operatorname{supp}(C)\cup
   (\operatorname{supp}(C)+\bar\omega)=\mathbb Z/D$.};
\end{tikzpicture}

\caption{The common-trace affine shift, the two labelled full targets, and the
exact odd-orbit inverse. The signed inverse shows why adding the two channels
does not create positivity from total coefficient mass.}
\label{fig:affine-channel}
\end{figure}

\section{Mixed-order sections with all integer coefficients retained}
Let $D\ge1$, $\lambda_i\in\mathbb Z$, and $N_i\ge1$. Write $C_x$ for the
coefficient of $[x]$ in~\eqref{eq:C}. Fix distinct coordinate indices
$i_1,\ldots,i_k$, and put
\[
 g_0=D,\qquad g_j=\gcd(g_{j-1},\lambda_{i_j}),\qquad
 m_j=g_{j-1}/g_j.
\]
The nontrivial chain has $g_0>g_1>\cdots>g_k=1$.

\begin{theorem}[Mixed-order coefficient completion]\label{thm:chain}
If $N_{i_j}\ge m_j$ for every $j$, then every original coefficient satisfies
\begin{equation}\label{eq:lower}
 C_x\ge L:=(\prod_{i\notin I}N_i)
              \prod_{j=1}^k\left\lfloor\frac{N_{i_j}}{m_j}\right\rfloor>0.
\end{equation}
The map
\begin{equation}\label{eq:radix}
 (t_{i_1},\ldots,t_{i_k})\longmapsto
 \sum_j\lambda_{i_j}t_{i_j}\pmod D,
 \qquad0\le t_{i_j}<m_j,
\end{equation}
is a bijection of the $D=\prod_jm_j$ original digit choices with $\mathbb Z/D$.
It has the explicit reverse algorithm below. When $N_{i_j}=m_j$ for all
selected coordinates, every coefficient equals $\prod_{i\notin I}N_i$.
The case $D=1$ is the empty chain, with every word full.
\end{theorem}
\begin{proof}
Let $H_g=g\mathbb Z/D\mathbb Z$. The quotient
$H_{g_j}/H_{g_{j-1}}$ has order $m_j$. Its specified generator is the image
of $\lambda_{i_j}$, since
$\gcd(g_{j-1}/g_j,\lambda_{i_j}/g_j)=1$.
For $c_k=x\pmod D$, compute backwards
\begin{equation}\label{eq:decode}
 t_{i_j}\equiv(c_j/g_j)(\lambda_{i_j}/g_j)^{-1}\pmod{m_j},\quad
 0\le t_{i_j}<m_j,\qquad
 c_{j-1}=c_j-\lambda_{i_j}t_{i_j}\pmod D.
\end{equation}
At step $j$, $c_j$ is divisible by $g_j$; the chosen digit makes the next
residual divisible by $g_{j-1}$. At the end it is zero modulo $D$.
The same quotient forces each digit of any other representation, proving
both existence and uniqueness.

Partition the selected interval $[0,N_{i_j})$ into its
$\lfloor N_{i_j}/m_j\rfloor$ full blocks of size $m_j$, with the remaining
short interval separately retained. For every choice of full blocks and every
choice of the unselected original coordinates, their offsets translate $x$.
Algorithm~\eqref{eq:decode} gives one unique original bounded word in those
blocks. Distinct choices give distinct words, proving~\eqref{eq:lower}.
When the selected intervals have exactly the radix lengths, this exhausts
all words and proves equality.
\end{proof}

\begin{theorem}[Partial-chain quotient reduction]
\label{thm:partial-chain}
Allow the same strict chain to end at a divisor $g=g_k$ of $D$, possibly
$g>1$, and retain $N_{i_j}\ge m_j$. Put
\[
 B=\prod_{j=1}^k\left\lfloor\frac{N_{i_j}}{m_j}\right\rfloor
\]
and, for $y\in\mathbb Z/g$, define the exact residual quotient coefficient
\[
 Q_y=\#\left\{(t_i)_{i\notin I}:0\le t_i<N_i,
 \quad\sum_{i\notin I}\lambda_it_i\equiv y\pmod g\right\}.
\]
Then there are integers $E_x\ge0$ such that
\begin{equation}\label{eq:partial-chain}
 C_x=BQ_{x\bmod g}+E_x.
\end{equation}
If $m_j\mid N_{i_j}$ for every selected coordinate, then every $E_x=0$.
For the two common channels of Theorem~\ref{thm:affine-channels}, with
$\delta=\omega\pmod g$,
\begin{equation}\label{eq:partial-pair}
 C_x+C_{x-\bar\omega}
 \ge B\bigl(Q_{x\bmod g}+Q_{(x\bmod g)-\delta}\bigr).
\end{equation}
Thus the complete-block subsource isolated by this partial chain forces
two-channel occupancy for every target precisely when
\begin{equation}\label{eq:quotient-cover}
 \operatorname{supp}(Q)+\{0,\delta\}=\mathbb Z/g.
\end{equation}
When $g>1$, write $d_g=\gcd(g,\delta)$. Since $g$ is odd,
condition~\eqref{eq:quotient-cover} requires the sharp bound
\begin{equation}\label{eq:quotient-support-bound}
 |\operatorname{supp}(Q)|\ge(g+d_g)/2.
\end{equation}
If
\[
 b=\gcd\bigl(g,\{\lambda_i:i\notin I\}\bigr)>1,
\]
with $b=g$ when no residual coordinate remains, then
condition~\eqref{eq:quotient-cover} is impossible.
\end{theorem}

\begin{proof}
For each selected coordinate, retain all complete radix blocks
\[
 a_jm_j+[0,m_j),\qquad
 0\le a_j<\lfloor N_{i_j}/m_j\rfloor,
\]
and retain its terminal interval separately. For fixed block indices the
offset is $\sum_j\lambda_{i_j}m_ja_j$. Since
\[
 \lambda_{i_j}m_j
 =g_{j-1}(\lambda_{i_j}/g_j),
\]
every such offset lies in $g\mathbb Z/D$. The bounded radix map on the
within-block digits is a bijection onto $g\mathbb Z/D$, by the same reverse
algorithm~\eqref{eq:decode} stopped at $g$. Therefore every unselected tuple
whose sum equals $x$ modulo $g$, together with every one of the $B$ block-index
tuples, supplies exactly one selected tuple whose total sum is $x$ modulo
$D$. These are $BQ_{x\bmod g}$ distinct original words. All words using at
least one terminal interval make up $E_x\ge0$. There is no terminal interval
when every $m_j$ divides $N_{i_j}$, proving equality in that case.

Apply~\eqref{eq:partial-chain} at $x$ and $x-\bar\omega$ to obtain
\eqref{eq:partial-pair}. Since all coefficients are nonnegative, its right
side is positive at every $x$ exactly under~\eqref{eq:quotient-cover}. On each
$\delta$-orbit in $\mathbb Z/g$, the complement of the support can have no
adjacent vertices. There are $d_g$ odd cycles of length $g/d_g$, giving
\eqref{eq:quotient-support-bound} by the odd-cycle argument already used in
Theorem~\ref{thm:affine-channels}; the same alternating construction attains
the bound.

Finally every residual sum in the complete-block subsource lies in the subgroup
$b\mathbb Z/g$. Its support
and one translate occupy at most two cosets of that subgroup. The index is the
odd integer $b\ge3$, so two cosets cannot cover $\mathbb Z/g$. This proves the
last assertion. In particular, the two channel labels cannot by themselves
replace a missing final radix: genuine residual quotient support is necessary.
Terminal intervals remain part of the original coefficient vector and may add
words outside this certificate.
\end{proof}

\paragraph{Carry, not a direct-product group substitution.}
The digit map is a set bijection. Its transported group operation is
$\operatorname{decode}(\operatorname{encode}(v)+\operatorname{encode}(w))$,
not coordinatewise addition in $\prod_j\mathbb Z/m_j$. For $D=9$ and steps
$(6,1)$ with radices $(3,3)$,
\[
 (0,2)\star(0,1)=(2,0),
\]
 because $3=2\cdot6\pmod9$. The cyclic extension and this integer carry remain.

\begin{figure}[ht]
\centering
\begin{tikzpicture}[
  font=\small,
  chain/.style={draw,rounded corners,minimum width=33mm,minimum height=8mm,align=center},
  every edge/.style={draw,-{Latex[length=2mm]},thick}
]
\node[chain] (h9) {$H_9=9\mathbb Z/9=\{0\}$};
\node[chain,right=16mm of h9] (h3) {$H_3=3\mathbb Z/9=\{0,3,6\}$};
\node[chain,right=16mm of h3] (h1) {$H_1=\mathbb Z/9$};
\path (h9) edge node[above] {$\lambda_1=6,\ m_1=3$} (h3)
      (h3) edge node[above] {$\lambda_2=1,\ m_2=3$} (h1);

\matrix[below=13mm of h3,matrix of math nodes,
  nodes={draw,minimum width=9mm,minimum height=7mm,anchor=center},
  row 1/.style={nodes={fill=black!6}},
  column 1/.style={nodes={fill=black!6}}] (grid) {
  {} & t_1=0 & t_1=1 & t_1=2 \\
  t_6=0 & 0 & 1 & 2 \\
  t_6=1 & 6 & 7 & 8 \\
  t_6=2 & 3 & 4 & 5 \\
};
\node[below=4mm of grid,draw,rounded corners,fill=black!4,inner sep=4pt]
  {$ (0,2)\star(0,1)=\operatorname{decode}(2+1)=(2,0)$};
\end{tikzpicture}

\caption{The exact $D=9$ subgroup chain and its carried digit section. The
entry in row $t_6$ and column $t_1$ is $6t_6+t_1\pmod9$. Every residue occurs
once, but the displayed carry proves that the section is not being replaced by
coordinatewise addition in $(\mathbb Z/3)^2$.}
\label{fig:mixed-chain}
\end{figure}

\paragraph{Fibres and coefficient loss.}
In the equality case each bin has
$\mu=\prod_{i\notin I}N_i$ words. Retaining the unselected digit vector and
applying \eqref{eq:decode} inverts the augmented map from a word to its residue
and that retained vector. The residue projection alone has $\mu$-point fibres.
Choose one reference word in every residue fibre. On free abelian modules, the
differences between each other word and that reference form a basis of the
kernel, of rank $D(\mu-1)=\mu D-D$. All pairwise differences generate the same
kernel but are not independent. The coefficient stabilizer is the whole
additive group.
If $\mu$ is zero modulo a prime, none of these original integer words vanishes.
Outside the equality case,~\eqref{eq:lower} counts a specified union of full
blocks; the remainder and its coefficients stay in~\eqref{eq:C}.

\begin{proposition}[Optimal bound within the radix-chain class]\label{prop:dag}
There is a finite exact algorithm for the largest lower bound~\eqref{eq:lower}
among all chains. Its graph has the divisors $g\mid D$ as vertices, and an edge
\[
 g\longrightarrow d=\gcd(g,\lambda_i)<g
 \quad\text{when}\quad N_i\ge g/d,
\]
with score $\lfloor N_i/(g/d)\rfloor/N_i$.
Start with score one at $D$, maximize products along decreasing paths, and
multiply the score at one by $\prod_iN_i$.
There are at most $n\tau(D)$ candidate edges. If one is unreachable, the zero
output means only that this certificate class fails, not that $C_{c_0}=0$.
\end{proposition}
\begin{proof}
After a coordinate has been used, every subsequent $g$ divides its step.
That coordinate can never give another strict edge. Thus paths automatically
have distinct labels. Future strict edges depend only on $g$, not on which
prefix reached it. The path score times $\prod_iN_i$ is exactly~\eqref{eq:lower}.
Dynamic programming on the decreasing divisor graph proves optimality within
this stated class. Exact rational scores become integers at every successful
terminal path by the displayed product formula.
\end{proof}

The criterion is complementary to, not a replacement for, the earlier
unit-width theorem. For $D=3$, steps $(1,1)$ and lengths $(2,2)$ give all three
bins, but neither length can start a radix chain. Conversely the examples
below pass mixed-order completion while their entire prior unit width is less
than $D-1$. Even joint generation and a positive source size are insufficient:
steps $(3,1)$ and lengths $(2,2)$ generate $\mathbb Z/9$ but miss five bins.

If an original full coefficient is absent, every feasible radix chain on its
trace part must therefore stop before one. The set of reachable divisors, and
the failed inequality $N_i<g/\gcd(g,\lambda_i)$ at each missing outgoing edge,
are an exact additional obstruction certificate. This does not assert that
such certificates are contradictory across all shells.

\section{Actual prime arithmetic with genuinely mixed orders}
The following are original bounded packets, not freely chosen residue supports.
Set $R=243$, $K=27$, $D=9$ and
\[
 a=H\cdot53\cdot163\cdot271,
 \qquad p=4a-243,
 \qquad H\equiv29\pmod{210}.
\]
For every prime member $p>R$ retain the M packet
\[
 \beta_{53}\in\{-1,1\},\quad
 \beta_{163},\beta_{271}\in\{-1,0,1\},
\]
with every other original centred exponent zero. If $H$ contains one of
these primes, the actual larger exponent interval is still retained and this
is its specified subpacket. Modulo $K$ the first prime is $-1$ and the other
two are one. The packet polynomial, up to its retained base shift, is
\[
 (1+X^5)(1+X^6+X^{12})(1+X+X^2)
 =2\sum_{c=0}^8X^c\pmod{X^9-1}.
\]
The chain is $9\to3\to1$ using steps $6,1$, both with radix three.
The preceding unit width is only $1+2=3<8$.
One full word is
\[
 u=H(163\cdot271)^2,\qquad
 (h,r,s,\lambda)=(H,44173,53,182),
\]
since $44173+53=243\cdot182$. Its original complement is the other full word.

There is an all-nonunit selected-chain example. Put
\[
 R=675,\quad K=45,\quad D=15,\quad
 a=H\cdot89\cdot271^2\cdot1801,\quad p=4a-675,
 \quad H\equiv101\pmod{210}.
\]
The original M packet has centred ranges $\{-1,1\}$ at $89$,
$[-2,2]$ at $271$, and $[-1,1]$ at $1801$, with the other coordinates zero.
Its polynomial is
\[
 (1+X^{11})\left(\sum_{j=0}^4X^{6j}\right)(1+X^{10}+X^{20})
 =2\sum_{c=0}^{14}X^c\pmod{X^{15}-1}.
\]
The selected steps $6,10$ are both nonunits; their chain is
$15\to3\to1$ with radices $5,3$. The only unit width is the single binary
$89$ coordinate, so the older unit-width criterion does not apply.
One full word and its marking are
\[
 u=H\cdot1801^2,\qquad
 (h,r,s,\lambda)=(H,1801,6536249,9686).
\]
Both examples are uniform original coefficient vectors with every bin equal
to two. Freezing either selected mixed direction at centred exponent zero
removes the full words of the specified M packet; the executable tables check
these exact restricted fibres as well.

These are infinite prime-source statements, not an assumption of prime values
of a nonlinear polynomial. With $A=\prod q_i^{e_i}$ fixed, $p=4AH-R$ and
$H=H_0+210n$ give a linear progression of modulus $840A$. It is reduced:
$p\equiv1\pmod{840}$ and $p\equiv-R\pmod{q_i}$ for each fixed prime $q_i$,
which is a unit. Dirichlet's theorem therefore supplies infinitely many prime
members \citep[Theorem~18.1]{Dirichlet}. Each such member has the constructed
source and full M return. This does not cover arbitrary prescribed primes.
The JSON tables select particular prime members by a bounded deterministic
search and certify their primality independently by complete-order certificates.

\subsection{Arithmetic sharpness at an arbitrary odd defect}
The last radix capacity cannot be lowered even on actual hard-prime packets.
Fix odd $D>1$ and integers $\lambda_1,\ldots,\lambda_k$ such that
\[
 g_0=D,\qquad g_j=\gcd(g_{j-1},\lambda_j),\qquad
 D=g_0>g_1>\cdots>g_k=1.
\]
Put $m_j=g_{j-1}/g_j$, let the last radix be
$m=m_k=g_{k-1}>1$, and put $R=3D^2$, $K=3D$. For $i<k$, choose distinct primes greater
than seven with
\[
 q_i\equiv1+K\lambda_i\pmod R.
\]
Choose a further prime
\[
 q_k\equiv-1+Kc\pmod R,\qquad 2c\equiv-\lambda_k\pmod D.
\]
Every displayed class is a unit modulo $R$, so Dirichlet supplies these primes.
For $i<k$ one has $\operatorname{ord}_K(q_i)=1$ and
$(q_i-1)/K\equiv\lambda_i\pmod D$. Since $K^2=3R$,
\[
 q_k^2=(-1+Kc)^2\equiv1-2Kc\pmod R,
\]
so $\operatorname{ord}_K(q_k)=2$ and
\[
 \frac{q_k^2-1}{K}\equiv-2c\equiv\lambda_k\pmod D.
\]
Take original exponents $e_i=(m_i-1)/2$ for $i<k$ and $e_k=m-1$. Let
$A=\prod q_i^{e_i}$, and impose
\[
 p\equiv1\pmod{840},\qquad p\equiv4A-R\pmod{4A^2}.
\]
The congruences agree modulo four and are reduced at every other prime.
Infinitely many prime members therefore have $a=(p+R)/4=AH$ with
$H\equiv1\pmod A$, so the stated prime exponents are exact.

Hold the other centred coordinates zero. For $i<k$ use all
$\beta_i\in[-e_i,e_i]$; for $q_k$ use the odd part
\[
 \beta_k\in\{-(m-2),-(m-4),\ldots,m-2\}.
\]
For $i<k$, square-zero expansion gives
$q_i^{\beta_i}\equiv1+K\lambda_i\beta_i\pmod R$; because $\beta_k$ is odd,
$q_k^{\beta_k}\equiv-1+Kc\beta_k\pmod R$. Hence every displayed word is an
original M trace and its original residue ratio and fine defect are
\[
 z=\prod_iq_i^{\beta_i}
 \equiv-1+K\left(c\beta_k-\sum_{i<k}\lambda_i\beta_i\right)\pmod R,
 \qquad
 d(\boldsymbol\beta)=c\beta_k-\sum_{i<k}\lambda_i\beta_i\pmod D.
\]
For $i<k$, the centred interval $[-e_i,e_i]$ has exactly $m_i$ members and is
a complete digit interval modulo $m_i$. The reverse decoder therefore maps
their product bijectively onto $H_m=m\mathbb Z/D$, whose order is
$D/m=\prod_{i<k}m_i$. Moreover $\gcd(c,m)=1$, because
$2c\equiv-\lambda_k\pmod m$, $m$ is odd, and $\gcd(\lambda_k,m)=1$. The
$m-1$ displayed values of $\beta_k$ are exactly the nonzero residue classes
modulo $m$, so their multiples by $c$ are the nonzero quotient classes of
$(\mathbb Z/D)/H_m$. Thus the actual fine-defect count
$\mu_x=\#\{u:d(u)=x\}$ is
\begin{equation}\label{eq:sharp}
 \mu_x=\begin{cases}0,&m\mid x,\\1,&m\nmid x.\end{cases}
\end{equation}
Its actual coefficient stabilizer is exactly $H_m$: a translation preserves
the zero set precisely when it lies in $H_m$. The source maps bijectively onto
the complement of $H_m$. In particular zero is absent, despite all preceding
capacities being exact.
This construction does not furnish or claim an ES counterexample: it proves an
exact empty subpacket while retaining every other original word.

The separate finite realization has
\[
 D=9,\quad K=27,\quad R=243,\quad
 (\lambda_1,\lambda_2)=(6,1),\quad(q_1,q_2)=(163,107),
\]
\[
 A=163\cdot107^2=1866187,\quad H=523,\quad
 a=976015801,\quad p=3904062961.
\]
Its six centred exponent pairs
$\beta_{163}\in\{-1,0,1\}$, $\beta_{107}\in\{-1,1\}$ have fine-defect vector
\[
 (\mu_0,\ldots,\mu_8)=(0,1,1,0,1,1,0,1,1),\qquad
 \operatorname{Stab}(\mu)=\{0,3,6\}=3\mathbb Z/9.
\]
This $H=523$ example is distinct from the displayed $H\equiv1\pmod A$ lift.
The reduced exact-exponent progression through it is
\[
 p=3904062961+2925429291933960t.
\]
Complete enumeration finds no full word in the displayed $R=243$ shell. The
same prime nevertheless has the exterior full witness
\[
 a_0=976015741,\quad R_0=3,\quad u_0=89,\quad
 (h,r,s,\kappa)=(89,1,10966469,1305009810),
\]
and therefore the exact E denominators
\[
 (976015741,1273710116719419210,453440801203675619490).
\]
Thus the boundary packet is not an ES counterexample. Complete prime, word,
progression, stabilizer and exterior-witness certificates are supplied.

\section{A factor exchange that permits controlled upward movement}
The supplied least-example theorem shows that a repair from \emph{every} trace
cannot always satisfy $a'\le a$ \citep[\S7]{Previous}. We retain another exact
integer invariant instead. For any original shell and any chosen divisor
$s\mid a$, write $n=a/s$ and define
\begin{equation}\label{eq:invariant}
 C=R+s,\qquad M=p+C=4a+s=(4n+1)s.
\end{equation}
Here $0<C<3p/2$ and $M<5p/2$.

\begin{theorem}[Complete labelled plus-factor fibre]\label{thm:plus}
At fixed $p,C$, every positive first-half integer cell with invariant
\eqref{eq:invariant} is exactly
\[
 s'\mid M=p+C,\quad M/s'\equiv1\pmod4,\quad0<C-s'<p,
\]
\begin{equation}\label{eq:cell}
 n'=(M/s'-1)/4,\quad R'=C-s',\quad
 a'=n's'=(M-s')/4.
\end{equation}
The map is inverted by retaining $s$ and $s'$. Its unit-ray exterior targets
are exactly those cells additionally satisfying $R'\mid M$. For each,
\begin{equation}\label{eq:plusreturn}
 (h',r',s',\kappa')=(n',1,s',M/R'-1),\qquad u'=n',
\end{equation}
with ordered denominators $(a',n's'\kappa',pn'\kappa')$.
The change of first denominator is exactly
\[
 a'-a=(s-s')/4.
\]
\end{theorem}
\begin{proof}
Equation~\eqref{eq:invariant} proves both directions of~\eqref{eq:cell}; its
congruence implies $R'\equiv-p\equiv3\pmod4$, so the strict range gives
$R'\ge3$ and $p/4<a'<p/2$. Because $s'\mid a'$ and $a'<p$, one has
$\gcd(p,s')=1$; moreover
\[
 \gcd(R',s')=\gcd(4n's'-p,s')=\gcd(p,s')=1.
\]
Since $M=(4n'+1)s'$, the original unit-ray E gate
$R'\mid4n'+1$ is therefore equivalent to $R'\mid M$. Equivalently,
$p+s'=M-R'$, so $R'\mid p+s'$ precisely when $R'\mid M$.
Then~\eqref{eq:plusreturn} is the inherited full E identity, and
$u'=n'$ divides $a'^2$ with its actual original exponents.
Conversely a full unit-ray target in this fibre forces that same divisibility.
Both residual and carrier factors are computed, not declared favourable.
\end{proof}

The whole cell list has at most $\tau(M)$ members. The full target list is a
specified finite divisor test, not an assertion that every fibre has a target.
For a canonical original E or M trace, choose $s$ from its exact gcd
factorization. Then $n=hr$. For a marked unit-ray E target, its complete
\emph{canonical} incoming trace fibre is: enumerate every cell
\eqref{eq:cell}, every $r\mid n$ with $\gcd(r,s)=1$, put
\[
 h=n/r,\qquad u=nr,
\]
and test each original square gate. This reconstructs all fields and all
original word bounds. An M word already retains its order. On ordered rational
solution triples define the exact E-tail involution
\[
 \tau_E(a,Y,Z)=(a,Z,Y),\qquad \tau_E^2=\operatorname{id}.
\]
This permutation does not generally preserve the original E exponent box, so it
is retained as a map of ordered solutions and is not inserted as another member
of the canonical incoming word fibre. In particular it is not an extra free
bit on an already counted M word.

The labelled arrow includes its chosen target. Forgetting that choice makes
the forward relation set-valued when more than one divisor passes. On free
integer coefficient modules, choose one reference arrow over every target. The
differences between each other arrow and its fibre reference form a kernel basis;
all pairwise same-target differences generate that kernel but are dependent.
No group-ring reduction changes the count
of distinct original source words or of distinct target states.

\subsection{The critical 67369 source moves up by one shell}
At
\[
 p=67369,\quad a=16849=7\cdot29\cdot83,\quad R=27,\quad u=4067,
\]
the exact M gcd factors are $(h,r,s)=(83,7,29)$ and $\lambda=4/3$.
Thus $C=56$, $M=67425$. Select
\[
 s'=25,\quad R'=31,\quad n'=674,\quad a'=16850.
\]
Both $25$ and $31$ divide $67425$, with $67425/25\equiv1\pmod4$.
The resulting marking is
\[
 (a',R',u',h',r',s',\kappa')=(16850,31,674,674,1,25,2174),
\]
and
\begin{equation}\label{eq:critical}
 \frac4{67369}=\frac1{16850}+\frac1{36631900}+\frac1{98714178844}.
\end{equation}
The old factors $7,29,83$ and the new factors $2,5,337$ do not share a prime.
The invariant is $4a+s=4a'+s'$, not an identification of their factor boxes.

\begin{figure}[ht]
\centering
\begin{tikzpicture}[
  font=\scriptsize,
  cell/.style={draw,rounded corners,align=center,minimum width=31mm,minimum height=18mm,inner sep=3pt},
  source/.style={cell,fill=blue!8,very thick},
  target/.style={cell,fill=green!10,very thick},
  fail/.style={cell,fill=black!4},
  arrow/.style={-{Latex[length=2.2mm]},very thick}
]
\node[font=\small] (head) {$p=67369,\quad C=56,\quad M=p+C=67425$};
\node[source,below=8mm of head,xshift=-54mm] (c1)
  {$s=29,\ R=27$\\$a=16849,\ n=581$\\M trace $u=4067$};
\node[target,right=5mm of c1] (c2)
  {$s'=25,\ R'=31$\\$a'=16850,\ n'=674$\\E full $u'=674$};
\node[fail,right=5mm of c2] (c3)
  {$s'=5,\ R'=51$\\$a'=16855,\ n'=3371$\\$51\nmid67425$};
\node[fail,right=5mm of c3] (c4)
  {$s'=1,\ R'=55$\\$a'=16856,\ n'=16856$\\$55\nmid67425$};
\draw[arrow] (c1) -- (c2);
\node[draw,rounded corners,below=6mm of c2,inner sep=4pt,align=center]
  {$\displaystyle\frac4{67369}=\frac1{16850}+\frac1{36631900}
       +\frac1{98714178844}$};
\end{tikzpicture}

\caption{The complete $p=67369$, $C=56$, $M=67425$ plus-factor cell set.
The arrow retains the source word and selected target. Only the $s'=25$ cell
passes the unit-ray terminal test $R'\mid M$; the other divisor cells remain in
the fibre rather than being discarded.}
\label{fig:plus-fibre}
\end{figure}

The release enumerates every divisor at residuals $3,7,11,15,19,23,27,31$.
The first seven have no full E/M word; $R=27$ has proper traces and $R=31$
has full E words $674,84250$. This rechecks the relevant fixed-prime obstruction;
it does not claim to have newly replayed the predecessor's leastness proof over
all 817 smaller primes. The chosen target's canonical incoming trace fibre is
computed by the complete rule above, including any full source itself.

The return extends to a reduced prime progression, not just one example. Put
\[
 h=83+4650k,\quad p=812h-27=67369+3775800k,\quad k\ge0.
\]
Every prime member is $169\pmod{840}$ and has the proper source
\[
 (a,R,u,h,r,s,\lambda)=(203h,27,49h,h,7,29,4/3).
\]
The target is
\[
 a'=203h+1,\quad R'=31,\quad u'=h'=(203h+1)/25,
 \quad r'=1,\quad s'=25,\quad\kappa'=(p+25)/31.
\]
All divisions are exact. The progression is reduced, so infinitely many
prime members exist. The old word and ray deletion systems fail on this
source: their moduli exceed $27$, while the properness requires deletion by
a multiple of three. The earlier factor-sum map to $\delta=3$ also fails
because its necessary divisibility is $7\mid(3^2-1)/4=2$.
The new formula always has $a'=a+1$. Only the initial prime is asserted to
have an entirely empty old full shell; no such extra assertion is made for
every prime in the progression.

\subsection{Why the new fibre is not automatically occupied}
At $p=97,a=40,u=2$ in M, the canonical $s$ is $20$, so $C=83$ and $M=180$.
The complete cell list is
\[
 (a',R',s')=(36,47,36),(40,63,20),(44,79,4).
\]
None has a unit-ray E target. Exhausting each of these three complete original
square-divisor boxes finds no full E or M state at all. The input is nevertheless
a proper M trace. This refutes universal occupancy of a fixed plus-factor fibre,
not ES or possible transformations between different such fibres.

\section{Remaining arithmetic obligation and verification scope}
The mixed-order theorem forces original target coefficients when its actual
quotient capacities pass; its inverse and all integer multiplicities are known.
The upward return crosses a previously certified obstruction to every
nonincreasing repair at its selected source. The two mechanisms are not
implicitly composable on every input: any application must preserve the shown
source, interval, channel and divisibility tests.

The affine-channel theorem now gives an exact alternative to a full radix
chain. A partial chain may stop at $g>1$ only if its surviving quotient support
and its $\delta$-translate cover $\mathbb Z/g$. Two channel labels do not
constitute the missing final radix: when the residual steps remain in a proper
subgroup, the support meets at most two of at least three quotient cosets.
Thus a proof through the two-channel route must construct the required quotient
support, or introduce a proved source-changing map that alters this obstruction.
Total word mass, a norm bound, or the injectivity of $I+T_{\bar\omega}$ cannot
replace that step.

If ES failed at a prescribed prime, every full coefficient in the complete
first-half source would be zero. Each trace part would fail both the previous
unit-width test and every feasible mixed-order chain. Every chosen plus-factor
fibre reachable by a proposed argument would have to miss its displayed
terminal divisor test. These are exact necessary conditions, not a contradiction
by definition. The remaining Stage~7 step is still to force at least one
original gate, with controlled changes between shells if required.

The standard-library replay exhausts the complete original boxes at primes
$p\equiv1\pmod{24}$ through its recorded bound and checks every squareful
trace part against literal gates. It additionally examines 161887 finite
additive polynomials, compares the divisor-DAG optimum with an independent
ordered-subsequence implementation, retains all nonzero carries, and verifies
complete original example boxes and inverse fibres. The separate checker imports
neither the main program nor predecessor code. Large example primes have full
order certificates whose factors are themselves certified; trial division is
used in discovery and as a separate finite check. Finite search caps, outcomes,
and file hashes are in the release receipt. A probable-prime label is not used
as a proof. There is no Lean build, independent human review, priority
conclusion, or new overall ES verification-range claim.

\subsection{Primality and finite-check boundary}
The isolated large primes have complete-order certificates. A nonterminal node
retains the full certified prime factorization of $n-1$ and a base $b$ with
\[
 b^{n-1}\equiv1\pmod n,\qquad
 \gcd(b^{(n-1)/q}-1,n)=1\quad(q\mid n-1).
\]
For any prime divisor $\ell$ of $n$, the order of $b$ modulo $\ell$ divides
$n-1$ and does not divide $(n-1)/q$ for any prime $q$ of $n-1$. Thus that order
is $n-1$, which divides $\ell-1$. It follows that $\ell\ge n$, proving $n$
prime. Terminal nodes below $10000$ are checked by complete trial division
through their integer square root. Discovery tests and completed certificates
are recorded separately. The arithmetic census tests every source word at its
finite bound; it is not an additional assertion about primes beyond that bound.
